\documentclass[12pt]{article}
\usepackage[a4paper,margin=2.5cm]{geometry}
\usepackage[T1]{fontenc}
\usepackage{lmodern}
\usepackage{xcolor}
\usepackage{qrcode}
\definecolor{tautan}{RGB}{0,70,160}
\usepackage[colorlinks=true,urlcolor=tautan,linkcolor=tautan]{hyperref}
\pagestyle{empty}

\begin{document}
\begin{center}

{\LARGE\bfseries Tautan Proyek}\\[4pt]
{\large Pemodelan Komputasi Sistem Kamera Digital}\\[2pt]
{\large Fisika Komputasi --- Kelompok 5}\\[1.5em]
\rule{\linewidth}{0.6pt}\\[2.5em]

% ---- GitHub ----
{\Large\bfseries Repositori GitHub \textnormal{(\textit{Source Code})}}\\[1.2em]
\qrcode[height=3.2cm]{https://github.com/SandyFauzi/Fisika-Komputasi-5}\\[1.2em]
\href{https://github.com/SandyFauzi/Fisika-Komputasi-5}%
{\large\texttt{github.com/SandyFauzi/Fisika-Komputasi-5}}\\[3em]

% ---- Google Drive ----
{\Large\bfseries Draft Makalah \textnormal{(Google Drive)}}\\[1.2em]
\qrcode[height=3.2cm]{https://drive.google.com/drive/folders/12tDzCv-M5J5ZluB_qR3SLH_I3gEVZk5c?usp=sharing}\\[1.2em]
\href{https://drive.google.com/drive/folders/12tDzCv-M5J5ZluB_qR3SLH_I3gEVZk5c?usp=sharing}%
{\large Buka folder draft makalah}

\vfill
{\small Pindai kode QR atau klik tautan di atas untuk mengakses berkas.}

\end{center}
\end{document}
